\scanpage{326}
\mspage{65}
% archive sheet number: 325
% ms. 1
effet
% ms. 2–3
\diagram{\begin{tikzcd}[ampersand replacement=\&,row sep=0pt,column sep=small]
a\arrow[r,"\gamma"] \& b\arrow[r,"\gamma'"] \& c\\
c\arrow[r,"{\gamma'^{\unsure{\circ|o}}}"] \& b\arrow[r,"\gamma^\circ"] \& a
\end{tikzcd}}
% ms. 4
la relation
% ms. 5
\[(\gamma'\unsure{\,.\,|\,}\gamma)^\circ=\gamma^\circ\gamma'^\circ,\tag*{$(*)$}\]
% ms. 6–9
ce qui est en accord avec l'intuition que l'opération $\gamma\longmapsto\gamma^\circ$ joue le rôle d'un passage à l'inverse (``à homotopie près''), et avec la formule
% ms. 9
\[(\gamma'\gamma)^{-1}=\gamma^{-1}\gamma'^{-1}\]
% ms. 10–11
pour deux flèches \deleted{\unsure{inv.}} $\gamma,\gamma'$ composables (dans l'ordre $\gamma',\gamma$).
% ms. between 11 and 12
\figure{\begin{tikzpicture}[x=1cm,y=1cm]\draw (0,.20) -- (.10,0);\end{tikzpicture}}

% ms. 12
Pour la propriété universelle de
% ms. 13
\[X\longrightarrow\Pi(X)\ \unsure{\sim|,}\]
% ms. 14
(qui permettra d'écrire
% ms. 15
\[\Pi(X)\simeq X\,\Sigma^{\!-1}\deleted{\unsure{\mathrm{Fl}}},\quad\text{avec }\Sigma=\mathrm{Fl}(X)),\]
% ms. 16–17
autant l'établir tout de suite dans le cadre général où on se donne une partie
% ms. 18
\[\Sigma\subset\mathrm{Fl}(X),\]
% ms. 19
et on définit, pour $\tau\in\mathrm{Ty}$,
% ms. 20
\[\underline{\mathrm{Ch}}_{\Sigma,\tau}(X)\subset\underline{\mathrm{Ch}}_\tau(X)\]
% ms. 21–22
comme la sous-catégorie pleine de $\underline{\mathrm{Ch}}_\tau(X)$ formée des chemins dans $\deleted{\unsure{\Sigma}}X$ dont \uline{\deleted{toutes} les flèches rétrogrades}
